\subsection{シンプソン公式}

台形公式では,小区間ごとに曲線を直線で置き換えた.直線の代わりに放物線で置き換えれば,さらに近似が良くなると期待できるだろう.放物線を1つ決めるには3点が要るので,区間の両端に加えて中点の値も使う.

中点を$m$,半分の幅を$h$とし,区間を$[m-h,m+h]$と書く.\,$y=x-m$と置いて中心を$0$へ移し,3点を通る二次式を
\[
 q(y)=A+By+Cy^2
\]
とする.\,$q(0)=f(m)$より$A=f(m)$である.また,\,$q(h)$と$q(-h)$を足せば$B$の項が消えて
\[
 f(m-h)+f(m+h)=2A+2Ch^2
\]
となるので
\[
 C=\frac{f(m-h)-2f(m)+f(m+h)}{2h^2}.
\]
積分でも$By$は奇関数として消える.したがって
\begin{align*}
 \int_{-h}^hq(y)\dd y
 &=2hA+\frac{2h^3}{3}C\\
 &=2hf(m)+\frac h3\{f(m-h)-2f(m)+f(m+h)\}\\
 &=\frac h3\{f(m-h)+4f(m)+f(m+h)\}
\end{align*}
を得る.これがシンプソン公式である.中点の係数$4$は,二次式の積分を計算した結果として決まるのだ.

この公式は三次式についても正確である.中心を移した三次式のうち,三次と一次の項は奇関数なので,積分も両端と中点の重み付き和も$0$になる.残った二次と定数の項については,放物線の積分そのものなので正確である.この「三次式まで誤差がない」という性質により,誤差は四階微分で評価できる.
